% Edited conversion of source pp. 34--35.
% The source photographs are reconstructed as scalable TikZ schematics.
\ReportHeading
  {Про устаткування в Інституті механіки НАН України для дослідження конструкцій спеціальної техніки під дією ударних хвиль у воді й повітрі}
  {В.~А. Максимюк, Є.~А. Сторожук}
  {Інститут механіки ім. С.~П. Тимошенка НАН України}
  {}

В Інституті механіки ім. С.~П. Тимошенка НАН України збереглося у робочому стані та використовується унікальне експериментальне устаткування, створене майже пів століття тому: вертикальний бак для дослідження дії ударних хвиль на конструкції у воді та ударна труба діафрагмового типу.

\begin{wrapfigure}{r}{0.36\textwidth}
  \vspace{-0.6\baselineskip}
  \centering
  \begin{tikzpicture}[x=0.52cm,y=0.52cm,line cap=round,line join=round]
    \draw[thick] (-2.5,0) ellipse (2.5 and 0.7);
    \draw[thick] (-2.5,0)--(-2.5,5.6);
    \draw[thick] (2.5,0)--(2.5,5.6);
    \draw[thick] (-2.5,5.6) arc (180:360:2.5 and 0.7);
    \draw[dashed] (2.5,5.6) arc (0:180:2.5 and 0.7);
    \draw (-2.45,3.2)--(2.45,3.2);
    \draw (-2.45,3.2) arc (180:360:2.45 and 0.58);
    \draw[dashed] (2.45,3.2) arc (0:180:2.45 and 0.58);
    \draw[thick] (-1.25,1.15) rectangle (1.25,2.65);
    \draw[thick] (-1.25,2.65)--(-0.75,3.15)--(1.75,3.15)--(1.25,2.65);
    \draw[thick] (-0.75,3.15)--(-0.75,1.65)--(-1.25,1.15);
    \draw[thick] (1.75,3.15)--(1.75,1.65)--(1.25,1.15);
    \draw[thin] (-1.5,5.9)--(-1.5,3.15);
    \draw[thin] (1.5,5.9)--(1.75,3.15);
    \node[font=\scriptsize] at (0,-0.2) {бак із водою};
  \end{tikzpicture}
  \caption{Схема вертикального бака.}
  \vspace{-0.5\baselineskip}
\end{wrapfigure}

Вертикальний бак заввишки 3~м і діаметром 3.2~м розташовано у заглибленому підвальному приміщенні. На середині висоти влаштовано прозоре вікно для спостережень. Водяна арматура та помпа забезпечують наповнення бака й відкачування води, а навколо верхньої частини розташовано огороджений робочий майданчик (рис.~1). Раніше ударну хвилю у воді формували дротиками, які вибухали внаслідок розряду високовольтних конденсаторів~\cite{equipment:anikiev1999}. Нещодавно бак використовували для визначення шару води, достатнього для сповільнення кулі до прийнятної швидкості.

Ударна труба завдовжки 4.75~м має прямокутний внутрішній переріз $140\times210$~мм і чистоту внутрішньої поверхні 5-го класу. Вона формує ступінчасту ударну хвилю з крутим фронтом і сталим надлишковим тиском до 100~кПа поза фронтом. Тривалість імпульсу становить приблизно 8~мс, а час наростання тиску у фронті --- близько 40~мкс. Установку застосовують для визначення залежності власних частот тонкостінних конструкцій від тиску, що дає змогу оцінювати їх стійкість без руйнування~\cite{equipment:anikiev2013}.

\begin{center}
\begin{tikzpicture}[x=0.82cm,y=0.82cm,line cap=round,line join=round]
  \draw[thick] (0,0) rectangle (11,1.25);
  \draw[thick] (0.35,0.18) rectangle (3.7,1.07);
  \draw[thick] (3.7,0.1)--(3.7,1.15);
  \draw[thick] (8.8,0.05) rectangle (10.1,1.2);
  \draw[thick] (10.1,0.2)--(11.55,0.2)--(11.55,1.05)--(10.1,1.05);
  \draw[thick] (1.0,0)--(0.75,-0.8)--(1.25,-0.8)--(1.5,0);
  \draw[thick] (5.2,0)--(4.95,-0.8)--(5.45,-0.8)--(5.7,0);
  \draw[thick] (9.6,0)--(9.35,-0.8)--(9.85,-0.8)--(10.1,0);
  \draw[thin] (5.4,1.25)--(5.4,2.0)--(6.9,2.0)--(6.9,1.25);
  \draw[thin] (7.45,1.25)--(7.45,2.35)--(8.3,2.35)--(8.3,1.25);
  \node[font=\scriptsize] at (5.5,-1.05) {ударна труба з вимірювальним обладнанням};
\end{tikzpicture}
\captionof{figure}{Схема ударної труби з допоміжним обладнанням.}
\end{center}

\renewcommand{\refname}{Література}
\begin{thebibliography}{9}
\bibitem{equipment:anikiev1999}
I.~I. Anik'ev, P.~Z. Lugovoi, and E.~A. Sushchenko,
``Experimental studies of the protective properties of flat shock-wave shields in water,''
\emph{International Applied Mechanics}, vol.~35, no.~3, pp.~312--314, 1999.

\bibitem{equipment:anikiev2013}
I.~I. Anik'ev, V.~A. Maksimyuk, M.~I. Mikhailova, and E.~A. Sushchenko,
``Nonstationary behaviour of a cantilever-rod system under nearly critical loads,''
\emph{International Applied Mechanics}, vol.~49, no.~5, pp.~570--576, 2013.
\end{thebibliography}
